\noindent\textbf{Por qué se usa este procedimiento.} La exigencia $k=EA/L$ fija un área distinta para cada material: $A=kL/E$. Luego $m=\rho AL=\rho kL^2/E$. Comparar densidad sola sería incorrecto porque el área cambia.

\begin{enumerate}

\item La masa es $m = \rho k L^2/E$.

\end{enumerate}

\noindent\textbf{Operaciones y resultados intermedios.}

\begin{enumerate}

\item Masa del tirante.

\[\fitmath{y_{1}=\frac{1600\cdot 50000000\cdot {2}^{2}}{1.81\times10^{11}}\approx 1.767956\,\mathrm{kg}}\]

\end{enumerate}
